\documentclass[11pt]{article}
\usepackage[margin=1in]{geometry}
\usepackage{amsmath,amssymb}
\usepackage{graphicx}
\usepackage{booktabs}
\usepackage{hyperref}
\usepackage{microtype}
\usepackage{parskip}
\usepackage{xcolor}
\usepackage{enumitem}

\setlength{\parskip}{6pt}
\setlength{\parindent}{0pt}

\title{\textbf{Gradient Surgery for Perturbed Optimizers}\\
\large Mathematical Motivation from a Probabilistic Viewpoint}
\author{Kyle Heuton}
\date{March 2026}

\begin{document}
\maketitle
\thispagestyle{empty}

\tableofcontents
\vspace{1em}
\hrule
\vspace{1em}

% ============================================================
\section{Introduction}
% ============================================================

The Differentiable Perturbed Optimizer (DPO; Berthet et al.\ 2020) promises to
train a prediction model by optimizing decision quality rather than prediction
accuracy.  In theory, a model that predicts costs well enough to make the right
combinatorial decision---even if its absolute predictions are biased---should
outperform one that minimizes squared error on the coefficients themselves.

In practice on the knapsack benchmark, the opposite is true.  DPO trained from
cold initialization achieves test regret $0.074$, while plain MSE regression
achieves $0.064$---the decision-focused method is \emph{worse} by 16\%.

Diagnostic experiments at the DPO local minimum reveal the mechanism:
\begin{itemize}[itemsep=2pt]
  \item The DPO gradient magnitude is large: $\|g_p\| = 15.1$ (not near zero).
  \item Yet 99.1\% of the gradient is orthogonal to the MSE improvement
        direction: $g_p \cdot g_m / (\|g_p\|\|g_m\|) \approx 0.009$.
  \item Per-instance DPO gradients are incoherent: cosine similarity uniform on
        $[-0.2,\, +0.4]$ across training instances.
  \item The prediction model has systematic negative bias $-1.6$ to $-4.2$ on
        all 20 items (vs.\ near-zero bias for MSE).
\end{itemize}

The DPO gradient is not useless---it encodes information about which decision
swaps are feasible under the current predictions---but it is mixed with a
dominant incoherent component driven by systematic prediction bias.  This document
motivates \textbf{gradient surgery}: at each training step, project out the
MSE-misaligned component of $g_p$ and inject the clean MSE gradient in its place.

Empirically, surgery achieves test regret $\mathbf{0.055}$, beating MSE (0.064),
$\lambda=10$ MSE regularization (0.061), and approaching the oracle ceiling
(0.052) within 0.003.

% ============================================================
\section{Problem Setup and Notation}
% ============================================================

We consider the standard Predict-then-Optimize (PtO) setting.  Let:
\begin{itemize}[itemsep=2pt]
  \item $(x_i, y_i) \in \mathcal{X} \times \mathbb{R}^D$ be $n$ training
        instances, where $x_i$ are features and $y_i$ are true cost coefficients.
  \item $\mathcal{Z} \subseteq \{0,1\}^D$ (or more generally $\mathbb{R}^D$) be
        the feasible set for the combinatorial optimizer.
  \item $z^*(y) = \arg\max_{z \in \mathcal{Z}}\, y \cdot z$ be the optimal
        decision under cost vector $y$.
  \item $f_\theta : \mathcal{X} \to \mathbb{R}^D$ be the prediction model with
        parameters $\theta$, and $\hat{y} = f_\theta(x) \in \mathbb{R}^D$ be
        the predicted costs.
  \item $R(\theta) = \mathbb{E}_{(x,y)}\!\left[y \cdot z^*(y) -
        y \cdot z^*(\hat{y})\right]$ be the expected decision regret.
\end{itemize}

The goal is to minimize $R(\theta)$ over $\theta$.  Note that $R(\theta) \geq 0$
with equality if and only if $z^*(\hat{y}) = z^*(y)$ almost surely.

For general CO problems, $z$ need not be binary and $f_\theta$ need not be
linear in $y$.  We will note where specific assumptions are invoked.

% ============================================================
\section{Two Views of ``Good Predictions''}
% ============================================================

\subsection{The MSE View: MLE under Gaussian Cost Noise}

If we assume $y_i = f_\theta(x_i) + \eta_i$ with $\eta_i \sim \mathcal{N}(0,
\sigma_y^2 I)$, the MLE estimator is:
\begin{equation}
  \theta_{\mathrm{MSE}} = \arg\min_\theta \sum_{i=1}^n
    \|y_i - f_\theta(x_i)\|^2.
  \label{eq:mse}
\end{equation}
The gradient w.r.t.\ predicted costs $\hat{y} = f_\theta(x)$ (before the chain
rule through $f_\theta$) is:
\begin{equation}
  g_m = 2(\hat{y} - y).
  \label{eq:gm}
\end{equation}
This is a \textbf{per-item bias corrector}: for each item $d$, $g_{m,d}$ pushes
$\hat{y}_d$ toward $y_d$ proportionally to the signed prediction error.  The
signal is coherent across all instances that share the same sign of error on
item $d$.  It sees the prediction error directly, not through the distorting
lens of the CO solver.

\subsection{The Decision View: What Information in $y$ Actually Matters?}

For a CO problem, only the \emph{ordering structure} of $y$ near the decision
boundary matters.  Specifically, the optimal knapsack solution changes only when
item values cross the current shadow price (the Lagrange multiplier on capacity).
Two cost vectors $y$ and $y'$ produce the same $z^*(y) = z^*(y')$ if they agree
on which items are above vs.\ below the shadow price.

This suggests that the Gaussian likelihood $\mathcal{N}(f_\theta, \sigma_y^2 I)$
assigns equal probability to cost vectors that may imply very different decisions.
A decision-focused likelihood would concentrate mass on cost vectors that yield
the correct decision.

One natural candidate is the Boltzmann distribution over feasible solutions:
\begin{equation}
  P(z^*(x) = z \mid x,\, \theta) \propto \exp\!\left(\beta \cdot
    f_\theta(x) \cdot z\right),
  \label{eq:boltzmann}
\end{equation}
with inverse temperature $\beta > 0$.  Maximizing the log-likelihood of
observed decisions under this model recovers the NCE and similar structured
prediction losses---not DPO.

\subsection{The DPO View: Expected Objective Under a Randomized Policy}

DPO takes a third approach.  Define the \emph{soft decision} under Gaussian
perturbations of the predicted costs:
\begin{equation}
  \bar{z}_\sigma(\theta, x) = \frac{1}{N}\sum_{n=1}^{N}
    z^*\!\left(f_\theta(x) + \sigma\varepsilon_n\right),
  \qquad \varepsilon_n \sim \mathcal{N}(0, I_D).
  \label{eq:zbar}
\end{equation}
The DPO loss is the regret of the soft decision under true costs:
\begin{equation}
  L_\sigma(\theta) = \mathbb{E}_{(x,y)}\!\left[
    \left|y \cdot \bar{z}_\sigma(\theta,x) - y \cdot z^*(y)\right|\right].
  \label{eq:dpo-loss}
\end{equation}

Interpretation: each $\varepsilon_n$ sample asks ``if the true costs were
$\hat{y} + \sigma\varepsilon_n$, what would we do?''  The average decision
$\bar{z}_\sigma$ is then evaluated under the actual true costs $y$.

\textbf{DPO backward pass (Berthet et al.).}
For Normal noise ($\nabla \log p(\varepsilon) = \varepsilon$), the Jacobian--vector
product through $\bar{z}_\sigma$ is:
\begin{equation}
  \frac{\partial L_\sigma}{\partial \hat{y}} = g_p
  = \frac{1}{N\sigma}\sum_{n=1}^{N}
    \left(z_n \cdot \frac{\partial L}{\partial \bar{z}}\right)\varepsilon_n,
  \label{eq:gp}
\end{equation}
where $z_n = z^*(\hat{y} + \sigma\varepsilon_n)$ and
$\partial L / \partial \bar{z}$ is the upstream gradient from the regret
objective through $\bar{z}$.

The gradient signal for coefficient $d$ is the noise sample $\varepsilon_{n,d}$,
weighted by how much $z_n$ aligns with the upstream gradient direction.
Crucially: a perturbation $\varepsilon_n$ contributes to $g_{p,d}$ \emph{only}
when it changes the decision---i.e., when $|\varepsilon_{n,d}|$ is large enough
to push $\hat{y}_d$ past the decision margin.

% ============================================================
\section{Why the DPO Gradient Becomes Incoherent}
% ============================================================

\subsection{Decision Switches Require Margin Crossings}

Let $m_d$ denote the \emph{decision margin} for item $d$: the amount by which
$\hat{y}_d$ must move to change item $d$'s membership in the optimal knapsack.
Formally, if $\lambda^*$ is the shadow price under $\hat{y}$, then
$m_d = |\hat{y}_d / w_d - \lambda^*|$ (value-per-weight minus shadow price).

A perturbation $\varepsilon_n$ changes the decision for item $d$ only if
$\sigma|\varepsilon_{n,d}| > m_d \cdot w_d$.  The gradient contribution for
item $d$ from instance $i$ is therefore:
\begin{equation}
  [g_p]_{i,d} \neq 0 \quad\Longleftrightarrow\quad
  \sigma|\varepsilon_{n,d}| > m_{i,d} \cdot w_d
  \;\text{ for at least one } n.
\end{equation}
Now suppose the prediction model has systematic bias: $\hat{y}_d \ll y_d$ for all
items $d$ (as we observe empirically---bias of $-1.6$ to $-4.2$).  Then:
\begin{enumerate}[itemsep=2pt]
  \item The shadow price $\lambda^*$ under $\hat{y}$ is set by different boundary
        items than under $y$.
  \item The ``swap pattern''---which items change membership under
        $\sigma\varepsilon_n$---reflects the margin structure of $\hat{y}$,
        not of $y$.
  \item Different instances have different \emph{predicted} margins, so each
        makes different swaps.  The per-instance gradients $g_{p,i}$ point in
        different directions.
  \item The average $g_p = (1/n)\sum_i g_{p,i}$ is nearly zero in any fixed
        direction---even though each $\|g_{p,i}\|$ is large.
\end{enumerate}
This is the mechanism: the DPO gradient is dominated by an incoherent
``which-boundary-items-to-swap'' signal that varies by instance, while the
coherent bias-correction signal (shared across all instances) is a small residual.

\subsection{Empirical Evidence on Knapsack}

At the DPO local minimum after 300 training epochs:
\begin{itemize}[itemsep=2pt]
  \item $\|g_p\| = 15.1$: the gradient is large, not near zero.
  \item $g_p \cdot g_m / (\|g_p\|\|g_m\|) \approx 0.009$: 99.1\% of $g_p$ is
        orthogonal to the MSE improvement direction.
  \item Per-instance cosine similarities between $g_{p,i}$ and $g_{p,j}$:
        uniform on $[-0.2,\, +0.4]$---no coherent signal at all.
  \item The prediction model has systematic negative bias on all 20 items:
        $\hat{y}_d - y_d \in [-4.2,\, -1.6]$ for the DPO model vs.\
        $[-0.1,\, +0.1]$ for the MSE model.
\end{itemize}

The bias is the root cause.  The DPO gradient never learned to correct the
absolute magnitude of cost predictions, only to make swaps in the distorted
space it operates in.  Paradoxically, the large $\|g_p\|$ is \emph{evidence} of
the problem: it reflects many incoherent swap signals canceling imperfectly,
not a strong coherent direction.

% ============================================================
\section{Gradient Surgery: Separating Decision Signal from Bias Noise}
% ============================================================

\subsection{Decomposing the DPO Gradient}

At each training step, we have two gradient signals w.r.t.\ the predicted
costs $\hat{y}$ (before the chain rule through $f_\theta$):
\begin{align}
  g_p &= \frac{1}{N\sigma}\sum_{n=1}^{N}
    \left(z_n \cdot \frac{\partial L}{\partial \bar{z}}\right)\varepsilon_n
  &&\text{[DPO gradient, Eq.~\eqref{eq:gp}]} \\[4pt]
  g_m &= 2(\hat{y} - y)
  &&\text{[MSE gradient, Eq.~\eqref{eq:gm}]}
\end{align}

Orthogonally decompose $g_p$ relative to $g_m$:
\begin{equation}
  g_p = g_p^{\parallel} + g_p^{\perp},
  \qquad
  g_p^{\parallel} = \frac{g_p \cdot g_m}{\|g_m\|^2}\, g_m,
  \qquad
  g_p^{\perp} = g_p - g_p^{\parallel}.
  \label{eq:decomp}
\end{equation}

\textbf{Interpretation of each component:}
\begin{itemize}[itemsep=2pt]
  \item $g_p^{\parallel}$: the component of the DPO gradient \emph{aligned with}
        the MSE direction.  This is DPO's implicit attempt to correct prediction
        bias---but it is high-variance and may have the wrong sign for individual
        instances (hence the scatter in per-instance cosines).  Empirically,
        $g_p^{\parallel} \approx 0.009\, g_p$ in magnitude.

  \item $g_p^{\perp}$: the component of the DPO gradient \emph{orthogonal to}
        MSE.  This is the unique information DPO provides that MSE does not:
        the decision-switch structure---which items are near the boundary under
        the current predictions.  This is the signal worth preserving.
\end{itemize}

\subsection{The Surgery Update}

The gradient surgery update replaces DPO's noisy bias-correction signal with
the clean MSE gradient while retaining the decision-switch signal:
\begin{equation}
  g_{\mathrm{surg}} = g_p^{\perp} + w \cdot g_m
    = \left(g_p - g_p^{\parallel}\right) + w \cdot g_m.
  \label{eq:surgery}
\end{equation}
The three terms have distinct roles:
\begin{enumerate}[itemsep=2pt]
  \item \textbf{Discard $g_p^{\parallel}$}: throw away DPO's noisy, possibly
        anti-aligned estimate of the bias direction.
  \item \textbf{Keep $g_p^{\perp}$}: retain the decision-switch information unique
        to DPO---which items are boundary candidates under the perturbed predictions.
  \item \textbf{Add $w \cdot g_m$}: inject the clean, low-variance MSE gradient
        for bias correction.
\end{enumerate}

\textbf{Parameter interpretation:}
\begin{itemize}[itemsep=2pt]
  \item $w = 0$: pure denoised DPO---removes anti-MSE corruption, keeps only the
        decision signal.
  \item $w \to \infty$: surgery collapses to pure MSE (decision signal overwhelmed).
  \item $w = 1$: ``replace the noisy $g_p^{\parallel}$ with the exact
        $g_m$''---the natural choice.
\end{itemize}

\textbf{Distinction from $\lambda$-MSE regularization.}
The $\lambda$-MSE combined loss adds $\lambda \cdot g_m$ directly on top of $g_p$:
\begin{equation}
  g_{\lambda\text{-MSE}} = g_p + \lambda \cdot g_m
    = g_p^{\parallel} + g_p^{\perp} + \lambda \cdot g_m.
  \label{eq:lam-mse}
\end{equation}
Comparing with Eq.~\eqref{eq:surgery}: if $g_p^{\parallel}$ is anti-aligned with
$g_m$ (which happens when DPO's bias-correction estimate has the wrong sign),
then $g_p^{\parallel}$ partially \emph{cancels} the beneficial $\lambda \cdot g_m$
term.  Surgery avoids this cancellation by first removing $g_p^{\parallel}$, then
adding the clean correction.  This is why surgery outperforms $\lambda$-MSE
regularization even at the optimal $\lambda = 10$.

\subsection{Probabilistic Interpretation}

Surgery can be understood as enforcing a \textbf{separation of concerns} that
matches the structure of the learning problem:

\begin{enumerate}[itemsep=4pt]
  \item \textbf{Systematic bias} (mean prediction error per item): this is a
        first-order moment problem.  $g_m$ is an unbiased estimator of
        $\nabla_\theta \mathbb{E}[\|y - \hat{y}\|^2]$ and directly corrects the
        mean error.  It is coherent across all instances: if $\hat{y}_d < y_d$ on
        average, $g_m$ pushes $\theta$ to increase $\hat{y}_d$ uniformly.

  \item \textbf{Decision-relevant ranking errors} (which items to swap at the
        boundary): this is an ordering problem.  $g_p^{\perp}$ captures variance
        in the decision-switch pattern across $N$ perturbation samples, orthogonal
        to the mean-correction direction.  It provides signal about the local
        structure of the feasible-set boundary that MSE cannot see.
\end{enumerate}

Viewed this way, surgery corresponds to a model where prediction errors have
two independent components: (a) an additive bias corrected by the MSE term,
and (b) a relative ranking noise corrected by the DPO term.  Surgery assigns
each component its natural gradient.

% ============================================================
\section{Assumptions and Generalization}
% ============================================================

\subsection{Does Surgery Require Linear CO Objectives?}

Berthet et al.'s convergence theory assumes $f(y, z) = y \cdot z$ (linear
objectives).  However, the \emph{backward pass formula} Eq.~\eqref{eq:gp} holds
for any CO problem where the solver is a pure $\arg\max$---the derivation requires
only that $z^*(\cdot)$ is well-defined and that $\varepsilon \mapsto z^*(\hat{y}
+ \sigma\varepsilon)$ is measurable.  Linearity in $y$ is not invoked in the
log-derivative trick.

Surgery operates entirely on the gradients $g_p$ and $g_m$ w.r.t.\ predicted
costs $\hat{y}$.  Neither requires linearity of $f$:
\begin{itemize}[itemsep=2pt]
  \item $g_p$ is the output of the DPO backward pass: valid for any CO problem,
        whether the objective is linear (knapsack, shortest path) or nonlinear
        (portfolio with a covariance term).
  \item $g_m = 2(\hat{y} - y)$: requires only that we have the true costs $y$.
        The formula follows from the MSE loss and does not reference the CO problem
        structure at all.
\end{itemize}

\textbf{Conclusion}: gradient surgery requires no additional assumptions beyond
DPO itself.  It applies to any PnO problem where predicted costs $\hat{y}$ are
directly comparable to true costs $y$ via squared error---knapsack, shortest
path, bipartite matching, portfolio, TSP with linear arc costs.

For non-additive objectives (e.g.\ stochastic programs where costs enter
nonlinearly), the MSE gradient $g_m$ still makes sense as a per-coefficient bias
corrector, even if it is not the gradient of the true generalization objective.
The surgery logic is unchanged.

\subsection{When Does Surgery Help?}

Surgery provides value when three conditions hold:
\begin{enumerate}[itemsep=2pt]
  \item \textbf{There is systematic bias}: $\mathbb{E}[\hat{y} - y] \neq 0$.
        If predictions are unbiased, $g_m$ has zero mean and adds noise rather
        than signal; surgery degenerates to denoised DPO ($w = 0$).

  \item \textbf{DPO gradient has anti-MSE components}: $g_p \cdot g_m < 0$ for
        a non-negligible fraction of steps or instances.  This is the condition
        under which $g_p^{\parallel}$ harms $\lambda$-MSE regularization.

  \item \textbf{The decision signal is present but noisy}: $g_p^{\perp} \neq 0$
        and carries information about the boundary structure.  If $g_p^{\perp}$
        were pure noise, removing $g_p^{\parallel}$ would leave us with useless
        noise plus $w \cdot g_m$---and surgery would degenerate to (possibly worse
        than) $\lambda$-MSE.
\end{enumerate}

Surgery is \emph{neutral} when $g_p \parallel g_m$ (DPO and MSE point the same
way): $g_p^{\perp} \approx 0$, and surgery reduces to
$g_{\mathrm{surg}} \approx w \cdot g_m$, which is equivalent to $\lambda$-MSE at
$\lambda = w$.

Surgery \emph{may hurt} only if $g_p^{\perp}$ is systematically misleading---but
since $g_p^{\perp}$ is averaged over $N = 100$ perturbation samples across the
batch, pathological behavior is unlikely.

\textbf{Practical heuristic}: surgery should be tried on any PnO problem where
MSE outperforms DPO from cold initialization, which is the common pattern across
problems in this benchmark.

% ============================================================
\section{Empirical Evidence}
% ============================================================

Results on the knapsack benchmark (20 items, 320 train / 80 val / 200 test
instances, DGP noise\_width$= 0.5$, DP solver with $N = 100$ perturbation samples,
300 training epochs, dense prediction model):

\begin{table}[h]
  \centering
  \caption{\textbf{Knapsack benchmark results.}
    Test regret (lower is better).  All perturbed runs use $\sigma = 0.5$.
    Surgery uses the DP solver; MSE uses Gurobi at test time.}
  \label{tab:results}
  \begin{tabular}{lcc}
    \toprule
    Method & LR & Test Regret \\
    \midrule
    Plain perturb ($\sigma = 0.5$)              & $5\text{e-}3$ & 0.074 \\
    MSE regularization $\lambda = 10$           & $5\text{e-}3$ & 0.061 \\
    MSE baseline (Gurobi)                       & $5\text{e-}2$ & 0.064 \\
    \textbf{Gradient surgery} $w = 0.5$         & $1\text{e-}2$ & \textbf{0.055} \\
    Oracle ceiling (noiseless DGP)              & ---           & 0.052 \\
    \bottomrule
  \end{tabular}
\end{table}

Key observations from Table~\ref{tab:results}:

\begin{itemize}[itemsep=4pt]
  \item \textbf{Surgery beats $\lambda$-MSE reg (0.055 vs.\ 0.061).}  Removing
        the anti-MSE component of $g_p$ before adding $g_m$ is strictly better
        than just adding $g_m$ on top.  This confirms that $g_p^{\parallel}$ was
        partially canceling the regularization effect.

  \item \textbf{Surgery beats MSE (0.055 vs.\ 0.064).}  The $g_p^{\perp}$
        component carries genuine decision-relevant signal that MSE cannot provide.

  \item \textbf{Surgery requires lr $= 1\text{e-}2$ (opposite of plain
        perturb's lr $= 5\text{e-}3$).}  Plain perturb needs a low LR to prevent
        gradient overshooting from the large incoherent $g_p$ component.  Surgery
        removes that corrupted component, allowing a larger effective step without
        instability.

  \item \textbf{cos\_sim($g_p$, $g_m$) $\approx 0.19$ throughout training.}
        Surgery helps even when $g_p$ and $g_m$ are weakly positively aligned.
        The projection fraction $g_p^{\parallel} / \|g_p\| \approx 0.19$ means
        roughly 19\% of the DPO gradient is being replaced with a cleaner signal
        at each step.

  \item \textbf{$w \in \{0.5, 1.0, 2.0\}$ gives similar results at lr $= 1\text{e-}2$
        (0.055--0.056).}  The surgery weight is not a sensitive hyperparameter,
        at least on this problem.
\end{itemize}

% ============================================================
\section{Limitations and Open Questions}
% ============================================================

We enumerate honest limitations with the current evidence.

\begin{enumerate}[itemsep=8pt]

  \item \textbf{Two forward passes.}  Surgery requires computing both the DPO
        loss and the MSE loss per step.  In our implementation, the DPO
        perturbation loop ($N = 100$ solver calls) dominates by $\sim 100\times$,
        so the MSE forward is cheap in practice.  But it is a non-zero overhead,
        and for solvers faster than DP it may matter.

  \item \textbf{New hyperparameter $w$.}  Surgery introduces the surgery weight
        $w$ alongside $\sigma$.  Empirically we found low sensitivity across
        $w \in \{0.5, 1.0, 2.0\}$ on knapsack, but this is a single problem and
        the joint $(w, \mathrm{lr})$ space has not been fully explored.

  \item \textbf{Is surgery just doing MSE?}  At $w \to \infty$, surgery collapses
        to MSE.  The fact that surgery at $w = 0.5$ outperforms $\lambda$-MSE
        regularization at the optimal $\lambda = 10$ (0.055 vs.\ 0.061) provides
        evidence that $g_p^{\perp}$ contributes genuine decision-focused signal,
        not merely that ``more MSE is better.''  But this should be tested on
        problems where MSE already achieves near-oracle performance---if surgery
        cannot improve over MSE there, it would suggest the $g_p^{\perp}$ benefit
        is limited to high-bias regimes.

  \item \textbf{No convergence theory.}  Berthet et al.'s theory applies to the
        unmodified DPO gradient.  Surgery changes the update direction in a way
        not covered by existing theory.  In particular, the surgery update
        Eq.~\eqref{eq:surgery} is not the gradient of any single scalar objective
        (it mixes the DPO gradient direction with the MSE gradient magnitude),
        so standard descent guarantees do not apply.  The empirical evidence is
        promising but the theoretical gap is real.

  \item \textbf{MSE auxiliary objective assumes squared loss.}  We use
        $g_m = \nabla_{\hat{y}} \|y - \hat{y}\|^2 = 2(\hat{y} - y)$, which
        implicitly assumes Gaussian cost noise.  For problems where costs are
        log-normal, bounded, or categorical, a different auxiliary objective may
        be more appropriate.  Surgery generalizes to any auxiliary gradient:
        replace $g_m$ with $\nabla L_{\mathrm{aux}}$ for whatever bias-correction
        loss is most natural for the problem.

  \item \textbf{Single problem, single seed.}  We validated on knapsack only
        (binary decisions, linear objective, Gaussian DGP, one random seed).
        The generalization argument in Section~\ref{sec:assumptions} is
        theoretical.  Empirical validation on portfolio (continuous decisions),
        shortest path (permutation structure), TSP, and bipartite matching is
        needed before claiming generality.

  \item \textbf{Learning rate interaction.}  Surgery \emph{requires} lr $=
        1\text{e-}2$, while plain perturb performs best at lr $= 5\text{e-}3$.
        This is consistent with surgery changing the effective gradient scale
        (removing the dominant high-variance component), not just the direction.
        The optimal lr may differ across problems, requiring joint tuning of
        $(w, \mathrm{lr})$ and potentially $\sigma$.

\end{enumerate}

% ============================================================
\section{Conclusion}
% ============================================================

Gradient surgery addresses the root cause of DPO underperformance on knapsack:
the DPO gradient conflates a noisy, sometimes anti-aligned bias-correction
component $g_p^{\parallel}$ with the genuine decision-switch signal $g_p^{\perp}$.
By projecting out $g_p^{\parallel}$ and replacing it with the clean MSE gradient,
surgery preserves the unique information DPO provides while eliminating the
incoherent noise from systematic prediction bias.

The method requires no additional assumptions beyond standard DPO, applies to any
CO problem with additive cost structure, and achieves test regret 0.055 on
knapsack---within 0.003 of the oracle ceiling (0.052) and better than both MSE
(0.064) and $\lambda = 10$ regularization (0.061).

The key open question is generalization: does the $g_p^{\perp}$ component carry
genuine decision-relevant information beyond knapsack?  The theoretical argument
(§\ref{sec:assumptions}) suggests yes whenever DPO's decision-switch signal is
meaningful.  Empirical validation on additional problems---particularly those where
the DPO gradient is less dominated by systematic bias---remains as future work.

% ============================================================
\appendix
\section{Connection to the Berthet et al.\ Backward Pass}
\label{app:berthet}
% ============================================================

For completeness, we trace through the Berthet et al.\ Jacobian--vector product
to make Eq.~\eqref{eq:gp} explicit.

The autograd function \texttt{perturbedSoftDecision} (from \texttt{perturbed.py})
computes:
\begin{align}
  \text{Forward:} &\quad \bar{z} = \frac{1}{N}\sum_{n=1}^N z_n, \quad
    z_n = z^*(\hat{y} + \sigma\varepsilon_n), \quad
    \varepsilon_n \sim \mathcal{N}(0, I_D). \\[4pt]
  \text{Backward:} &\quad g_d = \frac{1}{N\sigma}\sum_{n=1}^N
    \left(\sum_k [z_n]_k \cdot [\partial_{\bar{z}} L]_k\right) \cdot \varepsilon_{n,d}.
\end{align}

In code this is implemented via two einsum contractions:
\begin{enumerate}[itemsep=2pt]
  \item \textbf{Scores}: $s_n = z_n \cdot \partial_{\bar{z}} L$ \quad
        (shape: $(N, B)$)---scalar, ``how aligned is $z_n$ with the upstream gradient?''
  \item \textbf{Gradient}: $g_d = (1/N\sigma)\sum_n s_n \cdot \varepsilon_{n,d}$
        \quad (shape: $(B, D)$)---weighted average of noise samples.
\end{enumerate}

Item $d$ receives a large gradient $|g_{p,d}|$ when: (a) many samples $n$ have
$|s_n|$ large, meaning $z_n$ is far from the current soft decision, \emph{and}
(b) those samples have large $|\varepsilon_{n,d}|$, meaning item $d$'s cost was
heavily perturbed.  This is why the gradient is non-zero only when perturbations
are large enough to cause decision switches: $z_n = z^*(\hat{y} + \sigma\varepsilon_n)
\neq z^*(\hat{y}) = z_0$ only when $\sigma|\varepsilon_{n,d}|$ exceeds the decision
margin for some item $d$.

\end{document}
